\bibitem[1]{bib1} M. Akian, S. Gaubert, and A. Guterman. Tropical polyhedra are equivalent to mean payoff games. Int. J. Algebra Comput., 22(1):125001 (43 pages), 2012. \href{https://doi.org/10.1142/S0218196711006674}{doi:10.1142/S0218196711006674}.

\bibitem[2]{bib2} Marianne Akian, Ravindra Bapat, and Stéphane Gaubert. Max-plus algebra. In Chapter 25 in “Handbook of linear algebra”, volume 39 of Discrete Mathematics and Its Applications. Chapman and Hall, 2007.

\bibitem[3]{bib3} X. Allamigeon, P. Benchimol, S. Gaubert, and M. Joswig. Combinatorial simplex algorithms can solve mean payoff games. SIAM J. Optim., 24(4):2096–2117, 2014. \href{https://doi.org/10.1137/140953800}{doi:10.1137/140953800}.

\bibitem[4]{bib4} X. Allamigeon, P. Benchimol, S. Gaubert, and M. Joswig. Tropicalizing the simplex algorithm. SIAM J. Discrete Math., 29(2):751–795, 2015. \href{https://doi.org/10.1137/130936464}{doi:10.1137/130936464}.

\bibitem[5]{bib5} X. Allamigeon, P. Benchimol, S. Gaubert, and M. Joswig. What tropical geometry tells us about the complexity of linear programming. SIAM Rev., 63(1):123–164, 2021. \href{https://doi.org/10.1137/20m1380211}{doi:10.1137/20m1380211}.

\bibitem[6]{bib6} X. Allamigeon, S. Gaubert, and M. Skomra. Solving generic nonarchimedean semidefinite programs using stochastic game algorithms. J. Symbolic Comput., 85:25–54, 2018. \href{https://doi.org/10.1016/j.jsc.2017.07.002}{doi:10.1016/j.jsc.2017.07.002}.

\bibitem[7]{bib7} Eric Allender, Peter Bürgisser, Johan Kjeldgaard-Pedersen, and Peter Bro Miltersen. On the complexity of numerical analysis. SIAM J. Comput., 38(5):1987–2006, 2008/09. \href{https://doi.org/10.1137/070697926}{doi:10.1137/070697926}.

\bibitem[8]{bib8} D. Andersson and P. B. Miltersen. The complexity of solving stochastic games on graphs. In Proceedings of the 20th International Symposium on Algorithms and Computation (ISAAC), volume 5878 of Lecture Notes in Comput. Sci., pages 112–121. Springer, 2009. \href{https://doi.org/10.1007/978-3-642-10631-6_13}{doi:10.1007/978-3-642-10631-6\_13}.

\bibitem[9]{bib9} A. Atserias and E. Maneva. Mean-payoff games and the max-atom problem, 2009. URL: \href{https://www.cs.upc.edu/~atserias/papers/mean-payoff-games-max-atom/mpma.pdf}{https://www.cs.upc.edu/˜atserias/papers/mean-payoff-games-max-atom/mpma.pdf}.

\bibitem[10]{bib10} David Avis and Oliver Friedmann. An exponential lower bound for Cunningham’s rule. Math. Program., 161(1-2):271–305, 2017. \href{https://doi.org/10.1007/S10107-016-1008-4}{doi:10.1007/S10107-016-1008-4}.

\bibitem[11]{bib11} S. Basu, R. Pollack, and M.-F. Roy. Algorithms in Real Algebraic Geometry, volume 10 of Algorithms Comput. Math. Springer, Berlin, 2006. \href{https://doi.org/10.1007/3-540-33099-2}{doi:10.1007/3-540-33099-2}.

\bibitem[12]{bib12} A. Ben-Tal and A. Nemirovski. Lectures on Modern Convex Optimization: Analysis, Algorithms, and Engineering Applications, volume 2 of MOS-SIAM Ser. Optim. SIAM, Philadelphia, PA, 2001. \href{https://doi.org/10.1137/1.9780898718829}{doi:10.1137/1.9780898718829}.

\bibitem[13]{bib13} N. Bertrand, P. Bouyer-Decitre, N. Fijalkow, and M. Skomra. Stochastic games. In N. Fijalkow, editor, Games on Graphs. 2023.

\bibitem[14]{bib14} Marc Bezem, Robert Nieuwenhuis, and Enric Rodríguez-Carbonell. The max-atom problem and its relevance. In Proceedings of the International Conference on Logic for Programming, Artificial Intelligence, and Reasoning (LPAR), pages 47–61, 2008. \href{https://doi.org/10.1007/978-3-540-89439-1_4}{doi:10.1007/978-3-540-89439-1\_4}.

\bibitem[15]{bib15} Manuel Bodirsky, Mario Kummer, and Andreas Thom. Spectrahedral shadows and completely positive maps on real closed fields. J. Eur. Math. Soc., 2024.

\bibitem[16]{bib16} Manuel Bodirsky and Marcello Mamino. Tropically convex constraint satisfaction. Theory of Computing Systems, 62(3):481–509, 2018. An extended abstract of the paper appeared under the title “Max-Closed Semilinear Constraints” in the proceedings of CSR’16; preprint available under ArXiv:1506.04184.

\bibitem[17]{bib17} Endre Boros, Khaled Elbassioni, Vladimir Gurvich, and Kazuhisa Makino. A convex programming-based algorithm for mean payoff stochastic games with perfect information. Optim. Lett., 11(8):1499–1512, 2017. \href{https://doi.org/10.1007/s11590-017-1140-y}{doi:10.1007/s11590-017-1140-y}.

\bibitem[18]{bib18} Julian C. Bradfield and Igor Walukiewicz. The mu-calculus and model checking. In Edmund M. Clarke, Thomas A. Henzinger, Helmut Veith, and Roderick Bloem, editors, Handbook of Model Checking, pages 871–919. Springer, 2018. \href{https://doi.org/10.1007/978-3-319-10575-8_26}{doi:10.1007/978-3-319-10575-8\_26}.

\bibitem[19]{bib19} Cristian S. Calude, Sanjay Jain, Bakhadyr Khoussainov, Wei Li, and Frank Stephan. Deciding parity games in quasi-polynomial time. SIAM J. Comput., 51(2):17–152, 2022. \href{https://doi.org/10.1137/17M1145288}{doi:10.1137/17M1145288}.

\bibitem[20]{bib20} John Canny. Some algebraic and geometric computations in PSPACE. In Proceedings of the Twentieth Annual ACM Symposium on Theory of Computing (STOC), pages 460–467, New York, NY, USA, 1988. ACM. \href{https://doi.org/10.1145/62212.62257}{doi:10.1145/62212.62257}.

\bibitem[21]{bib21} Thomas Colcombet, Nathanaël Fijalkow, Pawel Gawrychowski, and Pierre Ohlmann. The theory of universal graphs for infinite duration games. Log. Methods Comput. Sci., 18(3), 2022. \href{https://doi.org/10.46298/LMCS-18(3:29)2022}{doi:10.46298/LMCS-18(3:29)2022}.

\bibitem[22]{bib22} A. Condon. The complexity of stochastic games. Inform. and Comput., 96(2):203–224, 1992. \href{https://doi.org/10.1016/0890-5401(92)90048-K}{doi:10.1016/0890-5401(92)90048-K}.

\bibitem[23]{bib23} Anne Condon. On algorithms for simple stochastic games. In Jin-Yi Cai, editor, Advances in Computational Complexity Theory, volume 13 of DIMACS Series in Discrete Mathematics and Theoretical Computer Science, pages 51–72. AMS, 1993.

\bibitem[24]{bib24} Yann Disser, Oliver Friedmann, and Alexander V. Hopp. An exponential lower bound for Zadeh’s pivot rule. Math. Program., 199(1):865–936, 2023. \href{https://doi.org/10.1007/S10107-022-01848-X}{doi:10.1007/S10107-022-01848-X}.

\bibitem[25]{bib25} Andrzej Ehrenfeucht and Jan Mycielski. Positional strategies for mean payoff games. International Journal of Game Theory, 8(2):109–113, 1979.

\bibitem[26]{bib26} Friedrich Eisenbrand, Matthieu Haeberle, and Neta Singer. An Improved Bound on Sums of Square Roots via the Subspace Theorem. In Wolfgang Mulzer and Jeff M. Phillips, editors, 40th International Symposium on Computational Geometry (SoCG 2024), volume 293 of Leibniz International Proceedings in Informatics (LIPIcs), pages 54:1–54:8, Dagstuhl, Germany, 2024. Schloss Dagstuhl – Leibniz-Zentrum für Informatik. \href{https://doi.org/10.4230/LIPIcs.SoCG.2024.54}{doi:10.4230/LIPIcs.SoCG.2024.54}.

\bibitem[27]{bib27} E. A. Emerson, C. S. Jutla, and A. P. Sistla. On model-checking fragments of $\mu$-calculus. In Proceedings of the 5th International Conference on Computer-Aided Verification (CAV), volume 697 of Lecture Notes in Comput. Sci., pages 385–396. Springer, 1993. \href{https://doi.org/10.1007/3-540-56922-7_32}{doi:10.1007/3-540-56922-7\_32}.

\bibitem[28]{bib28} E. Allen Emerson and Charanjit S. Jutla. Tree automata, mu-calculus and determinacy. In Proceedings of the 32nd annual Symposium on Foundations of Computer Science (FOCS’91), pages 368–377, Washington, DC, USA, 1991. IEEE Computer Society.

\bibitem[29]{bib29} John Fearnley, Spencer Gordon, Ruta Mehta, and Rahul Savani. Unique end of potential line. Journal of Computer and System Sciences, 114:1–35, 2020. \href{https://doi.org/10.1016/J.JCSS.2020.05.007}{doi:10.1016/J.JCSS.2020.05.007}.

\bibitem[30]{bib30} John Fearnley, Marcin Jurdziński, and Rahul Savani. Linear complementarity algorithms for infinite games. In SOFSEM 2010: Theory and practice of computer science. 36th conference on current trends in theory and practice of computer science, Špindlerův Mlýn, Czech Republic, January 23–29, 2010. Proceedings, pages 382–393. Berlin: Springer, 2010. \href{https://doi.org/10.1007/978-3-642-11266-9_32}{doi:10.1007/978-3-642-11266-9\_32}.

\bibitem[31]{bib31} O. Friedmann. A subexponential lower bound for Zadeh’s pivoting rule for solving linear programs and games. In Proceedings of the 15th International Conference on Integer Programming and Combinatorial Optimization (IPCO), volume 6655 of Lecture Notes in Comput. Sci., pages 192–206. Springer, 2011. \href{https://doi.org/10.1007/978-3-642-20807-2_16}{doi:10.1007/978-3-642-20807-2\_16}.

\bibitem[32]{bib32} Oliver Friedmann, Thomas Dueholm Hansen, and Uri Zwick. Subexponential lower bounds for randomized pivoting rules for the simplex algorithm. In Proceedings of the Forty-third Annual ACM Symposium on Theory of Computing (STOC), pages 283–292, New York, NY, USA, 2011. ACM. \href{https://doi.org/10.1145/1993636.1993675}{doi:10.1145/1993636.1993675}.

\bibitem[33]{bib33} Oliver Friedmann, Thomas Dueholm Hansen, and Uri Zwick. Errata for: A subexponential lower bound for the random facet algorithm for parity games. CoRR, abs/1410.7871, 2014. \href{https://arxiv.org/abs/1410.7871}{arXiv:1410.7871}.

\bibitem[34]{bib34} M. R. Garey, R. L. Graham, and D. S. Johnson. Some NP-complete geometric problems. In Eighth Annual ACM Symposium on Theory of Computing (Hershey, Pa., 1976), pages 10–22. Assoc. Comput. Mach., New York, 1976.

\bibitem[35]{bib35} Bernd Gärtner and Leo Rüst. Simple stochastic games and P-matrix generalized linear complementarity problems. In Fundamentals of computation theory. 15th international symposium, FCT 2005, Lübeck, Germany, August 17–20, 2005. Proceedings., pages 209–220. Berlin: Springer, 2005. \href{https://doi.org/10.1007/11537311}{doi:10.1007/11537311}.

\bibitem[36]{bib36} M. Grötschel, L. Lovász, and L. Schrijver. Geometric Algorithms and Combinatorial Optimization. Springer, Heidelberg, 1994. Second edition.

\bibitem[37]{bib37} Thomas Dueholm Hansen and Rasmus Ibsen-Jensen. The complexity of interior point methods for solving discounted turn-based stochastic games. In The nature of computation. Logic, algorithms, applications. 9th conference on computability in Europe, CiE 2013, Milan, Italy, July 1–5, 2013. Proceedings, pages 252–262. Berlin: Springer, 2013. \href{https://doi.org/10.1007/978-3-642-39053-1_29}{doi:10.1007/978-3-642-39053-1\_29}.

\bibitem[38]{bib38} Wilfrid Hodges. A shorter model theory. Cambridge University Press, Cambridge, 1997.

\bibitem[39]{bib39} Marcin Jurdzinski, Rémi Morvan, and K. S. Thejaswini. Universal algorithms for parity games and nested fixpoints. In Jean-François Raskin, Krishnendu Chatterjee, Laurent Doyen, and Rupak Majumdar, editors, Principles of Systems Design - Essays Dedicated to Thomas A. Henzinger on the Occasion of His 60th Birthday, volume 13660 of Lecture Notes in Computer Science, pages 252–271. Springer, 2022. \href{https://doi.org/10.1007/978-3-031-22337-2_12}{doi:10.1007/978-3-031-22337-2\_12}.

\bibitem[40]{bib40} Marcin Jurdziński and Rahul Savani. A simple P-matrix linear complementarity problem for discounted games. In Logic and theory of algorithms. 4th conference on computability in Europe, CiE 2008, Athens, Greece, June 15–20, 2008. Proceedings, pages 283–293. Berlin: Springer, 2008. \href{https://doi.org/10.1007/978-3-540-69407-6_32}{doi:10.1007/978-3-540-69407-6\_32}.

\bibitem[41]{bib41} L. Khachiyan. A polynomial algorithm in linear programming. Doklady Akademii Nauk SSSR, 244:1093–1097, 1979.

\bibitem[42]{bib42} Jan Křetínský, Emanuel Ramneantu, Alexander Slivinskiy, and Maximilian Weininger. Comparison of algorithms for simple stochastic games. Information and Computation, 289:104885, 2022. \href{https://doi.org/10.1016/J.IC.2022.104885}{doi:10.1016/J.IC.2022.104885}.

\bibitem[43]{bib43} J. B. Lasserre. Moments, Positive Polynomials and Their Applications. Ser. Optim. Appl. Imperial College Press, 2009. \href{https://doi.org/10.1142/p665}{doi:10.1142/p665}.

\bibitem[44]{bib44} Karoliina Lehtinen and Udi Boker. Register games. Log. Methods Comput. Sci., 16(2), 2020. \href{https://doi.org/10.23638/LMCS-16(2:6)2020}{doi:10.23638/LMCS-16(2:6)2020}.

\bibitem[45]{bib45} Karoliina Lehtinen, Pawel Parys, Sven Schewe, and Dominik Wojtczak. A recursive approach to solving parity games in quasipolynomial time. Log. Methods Comput. Sci., 18(1), 2022. \href{https://doi.org/10.46298/LMCS-18(1:8)2022}{doi:10.46298/LMCS-18(1:8)2022}.

\bibitem[46]{bib46} T. M. Liggett and S. A. Lippman. Stochastic games with perfect information and time average payoff. SIAM Rev., 11(4):604–607, 1969. \href{https://doi.org/10.1137/1011093}{doi:10.1137/1011093}.

\bibitem[47]{bib47} Donald A. Martin. Borel determinacy. Annals of Mathematics, 102(2):363–371, 1975. URL: \href{http://www.jstor.org/stable/1971035}{http://www.jstor.org/stable/1971035}.

\bibitem[48]{bib48} Rolf H. Möhring, Martin Skutella, and Frederik Stork. Scheduling with and/or precedence constraints. SIAM Journal on Computing, 33(2):393–415, 2004. \href{https://doi.org/10.1137/S009753970037727X}{doi:10.1137/S009753970037727X}.

\bibitem[49]{bib49} A.W. Mostowski. Games with forbidden positions. Tech. Rep. 78, University of Gdansk, 1991.

\bibitem[50]{bib50} Y. Nesterov and A. Nemirovskii. Interior-point polynomial algorithms in convex programming, volume 13 of Studies in Applied and Numerical Mathematics. SIAM, Philadelphia, PA, 1994. \href{https://doi.org/10.1137/1.9781611970791}{doi:10.1137/1.9781611970791}.

\bibitem[51]{bib51} Lorant Porkolab and Leonid Khachiyan. On the complexity of semidefinite programs. Journal of Global Optimization, 10(4):351–365, 1997. \href{https://doi.org/10.1023/A:1008203903341}{doi:10.1023/A:1008203903341}.

\bibitem[52]{bib52} Anuj Puri. Theory of hybrid systems and discrete event systems. University of California at Berkeley, 1995.

\bibitem[53]{bib53} M. L. Puterman. Markov Decision Processes: Discrete Stochastic Dynamic Programming. Wiley Ser. Probab. Stat. Wiley, Hoboken, NJ, 2005.

\bibitem[54]{bib54} Motakuri V. Ramana. An exact duality theory for semidefinite programming and its complexity implications. Mathematical Programming, 77:129–162, 1997. \href{https://doi.org/10.1007/BF02614433}{doi:10.1007/BF02614433}.

\bibitem[55]{bib55} C. Scheiderer. Spectrahedral shadows. SIAM J. Appl. Algebra Geom., 2(1):26–44, 2018. \href{https://doi.org/10.1137/17M1118981}{doi:10.1137/17M1118981}.

\bibitem[56]{bib56} Sven Schewe. From parity and payoff games to linear programming. In Proceedings of the 34th International Symposium on Mathematical Foundations of Computer Science (MFCS), pages 675–686, 2009. \href{https://doi.org/10.1007/978-3-642-03816-7_57}{doi:10.1007/978-3-642-03816-7\_57}.

\bibitem[57]{bib57} M. Skomra. Tropical spectrahedra: Application to semidefinite programming and mean payoff games. PhD thesis, Université Paris-Saclay, 2018. URL: \href{https://pastel.archives-ouvertes.fr/tel-01958741}{https://pastel.archives-ouvertes.fr/tel-01958741}.

\bibitem[58]{bib58} Ola Svensson and Sergei Vorobyov. Linear complementarity and P-matrices for stochastic games. In Perspectives of Systems Informatics. 6th International Andrei Ershov Memorial Conference, PSI 2006, volume 4378 of Lecture Notes in Comput. Sci., pages 409–423. Springer, 2006. \href{https://doi.org/10.1007/978-3-540-70881-0_35}{doi:10.1007/978-3-540-70881-0\_35}.

\bibitem[59]{bib59} Sergey P. Tarasov and Mikhail N. Vyalyi. Semidefinite programming and arithmetic circuit evaluation. Discrete Applied Mathematics, 156(11):2070–2078, 2008. In Memory of Leonid Khachiyan (1952-2005). \href{https://doi.org/10.1016/J.DAM.2007.04.023}{doi:10.1016/J.DAM.2007.04.023}.

\bibitem[60]{bib60} Uri Zwick and Mike Paterson. The complexity of mean payoff games on graphs. Theor. Comput. Sci., 158(1\&2):343–359, 1996. \href{https://doi.org/10.1016/0304-3975(95)00188-3}{doi:10.1016/0304-3975(95)00188-3}.

